%% Philosophy and the Mind Sciences article starter -- phimisci class.
%%
%% The phimisci class itself is NOT included in this download: it ships
%% with every modern TeX distribution (and with latextemplates.com's
%% compiler), so bundling a copy would only risk pinning you to an older
%% version than your own system already has.
%%
%% Class documentation: texdoc phimisci  |  https://ctan.org/pkg/phimisci
%% Licence: phimisci is LPPL-1.3c or later, (c) Ruhr-Universitaet Bochum.
%%
%% Check the journal's current author guidelines before submitting --
%% https://philosophymindscience.org -- since a bundled class can go
%% stale against a journal's live requirements.
%%
%% engine.txt: xelatex (required -- the class hard-errors with a
%% ClassError "wrong-tex-engine" under pdflatex; lualatex also works).
%%
%% NO BIBLIOGRAPHY IN THIS STARTER -- and that is deliberate, not an
%% oversight. This pipeline runs bibtex, never biber. phimisci loads
%% \RequirePackage[style=apa,natbib=true]{biblatex} with no explicit
%% backend= key, so biblatex's own default backend (biber) applies;
%% forcing \PassOptionsToPackage{backend=bibtex}{biblatex} was tried and
%% fails HARD at compile time with "Biber backend is required by style",
%% because biblatex-apa's apa.bbx calls \DeclareSortingTemplate, a
%% biber-only feature bibtex cannot process at all -- unlike lni/ijsra
%% elsewhere in this batch, there is no bibtex-compatible path here. If
%% your own workflow runs biber (most local TeX installs and Overleaf
%% do), delete this comment block, add back \addbibresource{your.bib}
%% and \printbibliography, and cite normally with \citep/\citet --
%% biblatex + biber + phimisci's apa style works fine, just not under
%% this site's bibtex-only compiler.
%% "stage" MUST be passed explicitly. \l__phimisci_output_authors_bool
%% starts life FALSE (\bool_new:N), and the class's `stage .default:n =
%% {preparation}` only fires if you write the bare key "stage" with no
%% value -- omitting the key entirely (as in a plain \documentclass
%% {phimisci} with no options) never runs any \__phimisci_stage_*: code
%% at all, so authors stay switched off and the header silently prints
%% "Anonymized" even though nothing about the document asked for blind
%% review. Confirmed by compiling without this option first: the author
%% line vanished from the title page with no warning or error.
%% stage=preparation shows authors+contact but suppresses the DOI/rights
%% footer (appropriate for a not-yet-submitted draft); use stage=final
%% once the paper is the journal's typeset, accepted version, or
%% stage=submission to reproduce anonymous-review formatting on purpose.
\documentclass[stage=preparation]{phimisci}

\usepackage{booktabs}

\title{Predictive Processing and the Boundaries of Cognitive Extension}

%% Authors are \and-separated within ONE \author{...} call, with
%% \affiliation{...} and (optionally) \orcid{...} inline per author.
\author{
  Ada Lovelace \affiliation{Institute of Analytical Engines, Berlin,
  Germany} \orcid{0000-0002-1825-0097}
  \and
  Second Author \affiliation{Department of Philosophy, Example
  University, USA}
}

\abstract{Predictive processing accounts of cognition are sometimes read
as offering a principled criterion for extended cognition: whatever
minimizes prediction error across an agent-environment coupling counts as
part of the cognitive system. We argue this reading conflates two
distinct notions of boundary -- functional and mechanistic -- and that
predictive processing, properly understood, is neutral on the extended
mind thesis rather than a vindication of it.}

\keywords{predictive processing; extended cognition; philosophy of mind;
active inference}

%% \contact{Name}{email} prints the "Primary contact:" line at the foot
%% of the title page -- omitting it leaves that label with nothing after
%% it, so set it whenever \l__phimisci_output_contact_bool is on (true
%% for every non-"submission" stage).
\contact{Ada Lovelace}{ada@example.edu}

\begin{document}

\maketitle

\section{Introduction}
The predictive processing framework has been enlisted on both sides of
the extended cognition debate. Proponents of extended cognition read the
framework's emphasis on agent-environment coupling as support for their
view; critics read the same framework as committed to a boundary drawn
firmly at the skull. We argue both readings overstate what the framework
itself entails.

\section{Two Notions of Boundary}
A \emph{functional} boundary is drawn wherever a process contributes to
minimizing prediction error; a \emph{mechanistic} boundary is drawn
wherever the physical implementation of that process actually resides.
Predictive processing, as a computational-level theory, speaks directly
only to the former.

\section{Why Extension Does Not Follow}
Table~\ref{tab:boundaries} summarises the two readings and the boundary
criterion each actually licenses.

\begin{table}[htb]
\centering
\caption{Two readings of predictive processing and their boundary
  criteria.}
\label{tab:boundaries}
\begin{tabular}{lll}
\toprule
Reading & Boundary criterion & Licenses extension? \\
\midrule
Extended-mind reading & Functional & Yes, if argued independently \\
Internalist reading   & Mechanistic (assumed) & No, but not from PP alone \\
\bottomrule
\end{tabular}
\end{table}

Neither reading follows from the predictive processing framework by
itself: functional continuity across a boundary is compatible with, but
does not entail, that the mechanisms realizing the relevant computations
extend beyond the organism.

\section{Conclusion}
Predictive processing is best read as neutral on the extended mind
thesis. Settling the extended cognition debate requires an independent
argument about mechanism, not a computational-level theory of prediction
error minimization alone.

\end{document}
